\documentclass[11pt,letterpaper]{article}
\PassOptionsToPackage{hyphens}{url}
\usepackage[margin=.83in,top=.8in,bottom=.8in,headheight=15pt]{geometry}
\usepackage[T1]{fontenc}
\usepackage{lmodern,amsmath,amssymb,mathtools,bm,booktabs,array,graphicx,tikz}
\usetikzlibrary{arrows.meta,positioning,calc,decorations.pathreplacing,patterns}
\definecolor{ink}{HTML}{17344D}\definecolor{accent}{HTML}{087E8B}\definecolor{pale}{HTML}{EDF5F5}\definecolor{gold}{HTML}{AD6B19}
\providecommand{\dd}{\,\mathrm{d}}\providecommand{\R}{\mathbb{R}}
\providecommand{\E}{\mathbb{E}}\providecommand{\Var}{\operatorname{Var}}
\providecommand{\Cov}{\operatorname{Cov}}\providecommand{\diag}{\operatorname{diag}}
\providecommand{\norm}[1]{\left\lVert#1\right\rVert}
\providecommand{\abs}[1]{\left|#1\right|}
\usepackage{microtype,hyperref,fancyhdr,enumitem,amsthm,sectsty,needspace}
\allsectionsfont{\raggedright}
\hypersetup{colorlinks=true,linkcolor=ink,urlcolor=accent,pdfauthor={Zan Ahmad},pdftitle={Rigid alignment and arclength representations of closed curves / Review}}
\setlength{\parindent}{0pt}\setlength{\parskip}{6pt}\linespread{1.07}
\setlist{itemsep=3pt,topsep=4pt,leftmargin=*}
\emergencystretch=2em\clubpenalty=10000\widowpenalty=10000
\pagestyle{fancy}\fancyhf{}\fancyhead[L]{\small Computational Cardiology}\fancyhead[R]{\small Lecture 13}
\fancyfoot[L]{\footnotesize Zan Ahmad}\fancyfoot[R]{\footnotesize\thepage}\renewcommand{\headrulewidth}{.3pt}
\newtheorem{proposition}{Proposition}\newtheorem{theorem}{Theorem}\newtheorem{lemma}{Lemma}
\theoremstyle{definition}\newtheorem{definition}{Definition}\newtheorem{example}{Example}
\numberwithin{equation}{section}
\begin{document}
\section*{Rigid alignment and arclength representations of closed curves --- review sheet}
Use this after attempting the worked examples and practice. Explain each claim in words before calculating.

\section*{Key ideas}
\begin{itemize}
\item Representation fixes what information is available.
\item RMS needs units and a correspondence rule.
\item Proper rigid alignment removes translation and rotation while retaining size and handedness.
\item Uniform polygonal arclength is a sampling convention with a fixed start and direction.
\item Implementation checks, geometric evidence and clinical validation answer different questions.
\end{itemize}
\section*{Common errors and repairs}
\textbf{Treating any equal-length arrays as corresponding.} Document how each index maps to the object.

\textbf{Calling the vertex mean a center of mass.} Specify the averaging measure and density before using that interpretation.

\textbf{Assuming resampling recovers missing geometry.} Interpolation stays on the supplied polygon.

\section*{Retrieve without your notes}
\begin{itemize}
\item What are the units of squared RMS and RMS?
\item Which determinant is required for a proper rotation?
\item Why is the closing endpoint not repeated?
\end{itemize}
\section*{Readiness check}
\begin{itemize}
\item Choose a representation and list the information it retains and discards.
\item Compute centering, a fixed-correspondence RMS and a proper rigid alignment diagnostic.
\item Explain why equal vertex counts do not establish correspondence.
\item Compare index matching with uniform polygonal arclength resampling and report remaining limitations.
\end{itemize}
\end{document}
